\documentclass{article}

% page layout
\usepackage[letterpaper, textwidth=125mm, textheight=230mm]{geometry}
\usepackage{setspace}
\setstretch{1.08}
\sloppy\sloppypar\raggedbottom\frenchspacing

\newcommand{\problemname}{Problem}
\newcounter{problem}
\newcommand{\problem}[1]{\refstepcounter{problem}\paragraph{\problemname~{\theproblem} --- {#1}:}}
\newcommand{\probref}[1]{\problemname~\ref{#1}}

\begin{document}
\section*{Observing strategy problem set}
\noindent
\textit{by} \textbf{David W. Hogg}, \textit{with input from \textbf{Niamh O'Sullivan} \& \textbf{Ryan Rubenzahl}.}

\paragraph{Introduction:}
The current observing plan (or cadence or strategy) for \textsl{Terra Hunting}
is that we will observe each star for one exposure each night.
There has been discussion of the exposure time, and a bit of push back on the every-star-once-every-night plan, but it remains the baseline.
I (Hogg) am convinced that observing each star once per night cannot be the correct strategy, for several reasons:
\begin{itemize}
\item
The every-star-once-every-night strategy is \emph{just made up}, not based on the optimization or quantitative comparison of possible strategies.
\item
The strategy fills out the final data set with many pairs of observations separated by integer multiples of the sidereal period, but very few pairs of observations that sample shorter intervals, and especially not intervals near 1/2, 3/2, 5/2, 7/2, \textsl{etc.} of sidereal.
That, in turn, probably reduces our sensitivity to short-period planets and noise sources with correlations on sub-day timescales.
\item
Single-exposure visits will probably be timed to minimze p-mode contributions to the final uncertainty budget.
It is almost always better to resolve and model nuisances than to hope that you can observe such that they don't matter.
Besides, long exposures bin the incoming photons, and binning is sinning.
\end{itemize}
This problem set is designed to find (and, we hope, solve) simple, quantitative problems that are extremely adjacent to the hard, real-world observing strategy questions that we face in 2026 and over the next decade.

It is Hogg's firmly held position that \emph{if} we can solve the problems in this problem set, then
\textsl{(a)}~we will have content for more than one publishable paper, and
\textsl{(b)}~we will have a good basis for making decisions about the \textsl{Terra Hunting} observing plan.

\problem{Information about a planet in an observing program}\label{prob:info}
There are many ways to detect or measure or characterize planets in radial-velocity data.
That said, there are strict bounds on how well we can possibly do, coming from information theory.
In particular, the Fisher information in a data set about the amplitude (say) of a signal with a certain period sets an absolute lower limit on the uncertainty that we can obtain by any unbiased method (see the Cram\'er--Rao bound).
In detail, there are many ways to search for and identify and characterize signals.
All these methods must deliver an uncertainty that is strictly not less than the uncertainty we compute by considering only the Fisher information on the amplitude of a sinusoidal signal at fixed period.

That is, we could search in period, amplitude, eccentricity, and orbital orientation space, and we could simultaneously marginalize out all noise sources and all nuisances, and we cannot do \emph{better} than the uncertainty implied by the Fisher information on the amplitude of a sinusoid at fixed period.

Now imagine that we have $n$ radial-velocity measurements $y_i$, taken at (exposure mid-point barycentric) times $t_i$, each of which as an uncertainty (assumed Gaussian and zero mean) $\sigma_i$.
What is the Fisher information on the amplitude $\kappa$ of a radial-velocity signal
\begin{equation}
y(t) = a\,\cos(\omega\,t) + b\,\sin(\omega\,t)
\end{equation}
where the angular frequency $\omega$ is related to the period $P$ by $\omega = 2\pi/P$,
and the amplitude $\kappa$ is related to the linear amplitudes $a,b$ by $\kappa^2=a^2+b^2$.

``What about eccentricity'' you ask?
Ignore it; if you can't measure the sinusoid, you definitely can't measure the eccentric Kepler curve.
Recall that we are obtaining a lower bound on any possible uncertainty.

In order to solve this problem, start by assuming that all of the observations are uncorrelated; that is, assume that the covariance matrix for the list of measurements $y_i$ is diagonal.
Then, what is the same answer for covariant measurements, and how do you parameterize that?

Choose a set of times for two cases:
In unrealistic case A the times are periodic at the sidereal period for $n=3650$ daily observations.
In unrealistic case B the $n=3650$ times are randomly chosen in a 10-year interval.
Plot the (optimistic, minimal, best-case) uncertainty on $\kappa$ as a function of period $T$.
You will have to choose some uncertainty or uncertainty distribution for your measurements; choose to match your best-case expectations for \textsl{Terra Hunting}.
More sophisticated noise models will be developed below.

\problem{Once per night}\label{prob:k}
Imagine that there are a fixed number of available hours available over the next 3650 nights, and imagine (incorrectly) that there is magically guaranteed availability for a particular magical target every single night.
Imagine (incorrectly) that every observation $i$ we make of this target gives us a measurement of the radial velocity $y_i$ of exactly the same uncertainty $\sigma_i=\sigma_0$, and that these uncertainties are well represented by zero-mean Gaussians of variance $\sigma_i^2$.
Most importantly, assume (incorrectly) that the observational uncertainties are all perfectly \emph{uncorrelated}:
Every observation's noise draw is independent from every other observation's.
That is, in this problem, we are assuming that the star has \emph{no non-trivial noise properties} and that the instrumentation is stable.

Now consider three possible programs:
\begin{itemize}
\item
In program A, the target is observed once per night.
\item
In program B, the target is observed a random number of times per night, where that number is drawn from a Poisson process with mean rate of unity.
\item
In program C, the target is observed $k$ times in a night, once every $k$ nights. Consider only small values of $k$, like $k\leq 6$.
\item
In program D, nights are chosen iid as ``observing nights'' with probability $1/k$.
On each chosen observing night, the target is observed $k$ times; on the not-chosen nights the target is not observed.
\end{itemize}
By construction, all of these programs use the same amount of integration time.
How does the sensitivity of the observing program to planetary companions at different periods depend on strategy?
You might have to make some additional decisions about how you distribute observations within a night.
We are happy if you simply assume that in the nights in programs B, C, and D with more than one observation, you literally do the intra-night observations one right after another (separated only by CCD read time).
(In this case it is worth noting that programs B, C, and D all use less wall-clock time at fixed integration time than program A, since they don't involve as much telescope slew and acquisition.)

This problem is extremely simplified but it will show (we think) that all four programs are equally sensitive (in a Fisher-information sense) to planets with orbital periods substantially longer than a few days.
It will also show (we think) that short-period signals are less detectable in program C than in programs A, B, and D,
but that otherwise programs A, B, and D will all be identical for the detection of all planets.

\problem{Visibility and weather in the decade}\label{prob:visibility}
Re-do \probref{prob:k} but having censored the four observing programs by visibility and weather expectations.
That is:
\begin{itemize}
\item
Choose a realistic celestial position (RA and Dec) for your target, and figure out on which nights it is visible according to some sensible airmass limits.
\item
Generate some random (and, if you are good, realistically correlated) bad weather, such that the target can only be observed on the non-bad nights.
\end{itemize}
Censor the four observing programs A, B, C, D you made in \probref{prob:k} such that they observe for fewer nights, but such that they are now actually realistic programs.
Check that all four programs still deliver approximately the same amount of total observation time.
How do your answers to \probref{prob:k} change in this more realistic case?
Obviously everything gets worse (in an information sense) as the total exposure time goes down;
more importantly, are there periods that become harder to detect?

\problem{Confirm Chaplin \textsl{et al.}}\label{prob:chaplin}
The Chaplin \textsl{et al.} paper shows plots of the RV uncertainty in a single exposure of time $T$,
given two sources of noise.
The first source is white noise, which is an unbeatable limit on our precision.
The second is noise from surface velocities of p modes (asteroseismic modes).
The paper starts with a power specrum for the noise and ends up with RV uncertainties.
Write code to do what Chaplin did, and reproduce their key figures.
Note that the Chaplin estimates are (extremely optimistic) Fisher information estimates, just like the ones we made in \probref{prob:info}.

When you execute this problem, you might be doing integrals over frequency of power spectra (and then taking square roots).
If you execute the problem this way, make sure you do those integrals sensibly.
One important issue: Should your integral include a point at \emph{zero frequency}?
That's a nasty thing: It corresponds to a random DC offset across the full data set.
Does that exist?
Our advice is to do any integrals you can analytically, analytically.
The white-noise component of the variance has straightforward analytic results.

Once that all works, and you can reproduce the Chaplin \textsl{et al.} results, figure out how to do a slightly different calculation:
Use the same Chaplin assumptions as you used in the first part of this problem, but use them to compute the \emph{covariance} between two data points, one starting at time $t_1$ with exposure time $T_1$, and the other starting at $t_2>(t_1 + T)$ with exposure time $T_2$.

\problem{Negatively correlated noise in a multi-exposre visit}\label{prob:negative}
The noise from p-mode oscillations is very different from white noise or red noise (like that from granulation, super-granulation, and spot rotation; see \probref{prob:gran}), because the p-mode oscillations create short-term \emph{negative} correlations in the noise or in the joint uncertainties of adjacent or near-in-time observations.
When you have a \emph{negative} noise correlation between two adjacent exposures, you can do \emph{better} than $\sqrt{2}$ when you combine them, just as, when you have a \emph{positive} correlation between two adjacent exposures, you do \emph{worse} than $\sqrt{2}$ when you combine them (because they share aspects of the noise).

In \probref{prob:chaplin}, you figured out how to compute variances and covariances for observations and pairs of observations, given both p-mode and white noise contributions.
Now imagine that you have $N$ observations, all of length $T$, with dead-time intervals (read times) $\Delta$ between them.
For this observing sequence, you can construct the $N\times N$ variance tensor you expect to represent the combined p-mode and white noise in this $N$-exposure sequence, treated as one ``visit.''

We want to consider the situation in which the observation is modeled as a multivariate ($N$-dimensional) Gaussian draw from a Gaussian with a mean $N$-vector and a $N\times N$ variance tensor.
For this problem we will imagine that the star has, over the visit, a fixed mean velocity $\mu$ (which is what we are trying to measure).
Thus the mean $N$-vector is treated as just being the mean $\mu$ repeated $N$ times.
The variance tensor is built as just described.

The first question to answer is: What is the Fisher information you obtain about the mean $\mu$ from one of these $N$-exposure visits?
It looks like the sum of all the $N^2$ terms in the inverse of the variance matrix!

The second question is: What is the (optimistic) measurement uncertainty $\sigma$, and (more importantly) Fisher information per wall-clock time, coming from different realistic scenarios?
Consider exposure times in the range 30-1500 seconds, and $1\leq N\leq 10$.
In computing wall-clock time, set the read-noise time to be reasonable and include a realistic slew time at the end (or beginning if you prefer).

In your answer to this, you will see Chaplin \textsl{et al.} as the $N=1$ results,
you will see that multi-exposure visits don't just do better, but do better \emph[per unit wall clock time],
and you will see that the best exposure time at $N=1$ is the worst exposure time at $N=10$.
Enjoy!

Explain why the only metric we could possibly care about is information per wall-clock time.
Okay fine, this is the only metric we could care about \emph{when the assumptions of this problem are actually correct} (and they aren't).

Note, in all of this, the inverses of variance tensors appear.
It might sound counter-intuitive, but you should never use your linear-algebra package's \texttt{inv()} function, for reasons of numerical stability.
Always use \texttt{solve()} and never \texttt{inv()}.
If you have trouble figuring out how to get rid of your \texttt{inv()} calls, get in touch with Hogg for advice.

\problem{Weather within the visit}
Now imagine that the transparency varies \emph{within} an exposure.
In \probref{prob:chaplin} you (probably) did some integrals of the power spectrum times sinc functions.
These sinc functions appear because the sinc is the fourier transform of the tophat.
If the transparency varies, the exposure function or window function for the observation is no longer a top hat,
but instead has some non-trivial structure in time.
Add non-trivial structure, by (say) modulating the transparency during each observation with a random polynomial (make sure it is always bounded in $[0,1]$) or a thresholded or sigmoided polynomial or something like that.
Show that the Chaplin \textsl{et al.} result is very sensitive to this variation in transparency.
After all, the special exposure times they find depend on zeroing out p-modes.
Transparency variations will break that ``resonance.''

Show that the information obtained in the multi-exposure visits computed in \probref{prob:negative} is much less sensitive to such variations.
To perform this analysis, you will have to assume that we \emph{know} how the transparency varied during the exposure.
Make this assumption, and, when we are in operations, we will find out whether or not that is substantially valid.
For extra credit, show that multi-exposure visits generally have sensitivity (in an inverse variance sense) that varies linearly with the total average transparency during the exposures.

\problem{Granulation and super-granulation}\label{prob:gran}
The most important issue for Earth hunting is almost certainly the correlated noise from granulation, super-granulation, and rotating surface features like spots.
These noise sources introduce \emph{positive} correlations over hours to days, and positive correlations are ruinous:
When two observations are positively correlated, there is no way to combine them such that you get a measurement
that is $\sqrt{2}$ better than each individual measurement.
Thus you don't get twice as much information when you spend twice as much time.

Model a reasonable granulation and super-granulation noise source with NIAMH WHAT KERNEL FUNCTION OR POWER SPECTRUM? WHAT AMPLITUDE?
Now repeat all of the calculations of \probref{prob:visibility} but including this term in the covariance matrix.
That is ignore (for now) p-modes, but include realistic white noise (individual-exposure unbeatable uncertainty) plus the granulation and super-granulation noise contributions.
Do the visibility-and-weather-censored observing plans still fare similarly, or do plans C and D take the lead for any values of $k$?

Note that this calculation might be technically hard, because you will have to evaluate a Gaussian Process on thousands of data points.
If this is prohibitive, and you don't want to learn anything about executing GPs, you can do the following hack:
Each observing season (from when the star first becomes visible in (say) the Spring to when the star is last visible in (say) the Fall) can be treated individually, and the information from each season (the Fisher information or the inverse variance) can be coadded over seasons.
Each individual season should be a tractable GP evaluation even with brain-dead code.

\problem{Maximizing information about the noise}\label{prob:noise}
There are combinatoric options for observing programs.
We can come up with more possible observing programs than there are atoms in the visible universe.
Therefore we won't try to truly optimize over all possible observing programs.
However, we can compare the four visibility (and weather) censored observing programs from \probref{prob:visibility} on various grounds.
We have compared them on detectability of planets as a function of planet period, and we will (below) compare them on mitigation of the (effective) noise coming from short-period planets.

Here we ask: Which, of (the visibility and weather censored versions of) observing plans A, B, C, D is most sensitive to the properties of the ``red noise'' from granulation, super-granulation, and spots.
HOGG FORMULATE THE QUESTION SPECIFICALLY.

HOGG NOTE: We will have to consider scheduling \emph{within the night}, and not just scheduling at the night level, if we are going to get this right.

\problem{Forest of planets}\label{prob:planets}
HOGG WRITE THIS...

Note that there are many ways to structure your code here, but in many of them, you will find that you need to ``take the inverse'' of the same covariance matrix over and over again.
You might be tempted to use \texttt{inv()} and then repeatedly re-use that inverse tensor, but see comments at the end of \probref{prob:negative}, where we said that you should always use \texttt{solve()} instead.
If you need to do multiple \texttt{solve()} operations with the same matrix, perform a Cholesky factorization, save the factorization, and then use \texttt{cho\_solve()}.
That's just as fast as the dumber solution, and way more numerically stable.

\problem{What if we can mitigate granulation?}
foo and bar

\problem{Putting it all together}
foo and bar.

\end{document}
